\IfFileExists{../main.tex}{%
  \documentclass[../main.tex]{subfiles}%
}{%
  \documentclass[main.tex]{subfiles}%
}
\begin{document}
\setcounter{chapter}{6}

% ===========================================================================
\section{Kết luận}
\label{section:6.1}

Đồ án đã xây dựng hoàn chỉnh một \textbf{hệ thống kết nối chia sẻ thực phẩm dư thừa} cho bối cảnh Việt Nam, gồm ba sản phẩm phối hợp: ứng dụng di động (React Native/Expo) cho ba vai trò người dùng, trang quản trị web (React) cho quản trị viên, và máy chủ phụ trợ (Node.js/Express + MongoDB) cung cấp REST API kết hợp kênh thời gian thực. Hệ thống hiện thực được trọn vẹn bài toán đặt ra ở Chương~1: rút ngắn khoảng cách thông tin giữa \textit{người có thực phẩm dư}, \textit{người cần nhận} và \textit{lực lượng vận chuyển trung gian}, theo thời gian thực và trong phạm vi địa lý gần.

\subsection{So sánh với các sản phẩm tương tự}
\label{subsection:6.1.1}

So với các giải pháp đã khảo sát ở Chương~1, sản phẩm của đồ án là giải pháp duy nhất đồng thời thoả mãn cả năm tiêu chí quan trọng (Bảng~\ref{tab:ss_sanpham}): miễn phí cho người nhận, hỗ trợ vai trò tình nguyện viên có điều phối tự động, kết nối thời gian thực, cho phép người nhận \textit{chủ động} đăng yêu cầu (luồng hai chiều), và phù hợp thị trường Việt Nam (hỗ trợ tiếng Việt).

\begin{table}[H]
\centering
\renewcommand{\arraystretch}{1.3}
\resizebox{\textwidth}{!}{%
\begin{tabular}{|l|c|c|c|c|}
\hline
\textbf{Tiêu chí} & \textbf{Too Good To Go} & \textbf{OLIO} & \textbf{Foodbank} & \textbf{Sản phẩm đồ án} \\ \hline
Miễn phí cho người nhận            & Không & Có    & Có          & \textbf{Có} \\ \hline
Hỗ trợ tình nguyện viên            & Không & Không & Có (thủ công) & \textbf{Có (tự động)} \\ \hline
Kết nối thời gian thực             & Có    & Có    & Không       & \textbf{Có} \\ \hline
Người nhận chủ động đăng yêu cầu   & Không & Không & Không       & \textbf{Có} \\ \hline
Phù hợp / hỗ trợ tiếng Việt        & Không & Không & Có          & \textbf{Có} \\ \hline
\end{tabular}%
}
\caption{So sánh sản phẩm đồ án với các giải pháp hiện có}
\label{tab:ss_sanpham}
\end{table}

Khác biệt cốt lõi nằm ở chỗ sản phẩm triển khai \textbf{đầy đủ mô hình ba vai trò} với cơ chế điều phối tự động (người cho -- người nhận -- tình nguyện viên), hỗ trợ \textbf{luồng hai chiều} (cả người cho đăng đơn lẫn người nhận đăng yêu cầu), điều mà các nền tảng quốc tế (Too Good To Go, OLIO) lẫn mô hình thủ công trong nước (Foodbank) đều chưa làm trọn vẹn.

\subsection{Kết quả đạt được}
\label{subsection:6.1.2}

Trong quá trình thực hiện, đồ án đã hoàn thành các nội dung chính sau:
\begin{itemize}
  \item \textbf{Quản lý tài khoản \& phân quyền:} đăng ký, xác thực OTP, đăng nhập, khôi phục mật khẩu; bốn vai trò (Donor, Receiver, Volunteer, Admin) với kiểm soát truy cập theo vai trò (RBAC).
  \item \textbf{Nghiệp vụ chia sẻ hai chiều:} người cho đăng đơn quyên góp kèm ảnh/loại/số lượng/hạn dùng; người nhận tìm theo khoảng cách và kết nối, hoặc chủ động đăng yêu cầu để người cho đáp ứng.
  \item \textbf{Điều phối giao nhận:} hai phương thức tự lấy / ủy thác tình nguyện viên; điều phối tình nguyện viên lân cận theo cơ chế ``ai nhận trước thắng''; theo dõi trạng thái và xác nhận bằng mã lấy hàng.
  \item \textbf{Tương tác thời gian thực:} trò chuyện 1--1 trong phạm vi đơn (Socket.IO), thông báo đẩy và in-app đa ngôn ngữ.
  \item \textbf{Ghi nhận đóng góp:} điểm thưởng, thống kê cá nhân, đánh giá hai chiều.
  \item \textbf{Quản trị:} trang web cho quản trị viên giám sát người dùng, đơn và phản hồi.
\end{itemize}
Các \textbf{đóng góp kỹ thuật nổi bật} (ghép cặp an toàn khi tương tranh, ước lượng khoảng cách lai có dự phòng, máy trạng thái tự phục hồi, chống gian lận giao nhận, xác thực thu hồi tức thì, thông báo hợp nhất) đã được trình bày chi tiết ở Chương~5.

\subsection{Hạn chế còn tồn tại}
\label{subsection:6.1.3}

Bên cạnh kết quả đạt được, đồ án vẫn còn một số hạn chế cần thẳng thắn nhìn nhận:
\begin{itemize}
  \item \textbf{Gửi OTP qua SMS chưa được tích hợp:} cơ chế OTP (sinh, băm, hết hạn, giới hạn số lần) đã hoàn thiện, nhưng việc gửi mã tới điện thoại hiện mới ở mức môi trường phát triển; cần tích hợp nhà cung cấp SMS để dùng thực tế.
  \item \textbf{Mới triển khai trên Android:} nền tảng Expo cho phép mở rộng sang iOS nhưng phạm vi đồ án chưa build và kiểm thử bản iOS.
  \item \textbf{Lưu trữ tệp cục bộ:} ảnh đại diện/món ăn/chat đang lưu trên máy chủ và phục vụ qua đường dẫn tĩnh, chưa dùng dịch vụ lưu trữ đối tượng (object storage)/CDN nên hạn chế khi mở rộng.
  \item \textbf{Chưa có bộ kiểm thử tự động:} việc kiểm thử chủ yếu thủ công/hộp đen; chưa có unit/integration test và quy trình CI/CD.
  \item \textbf{Truy vấn vị trí theo quét tuyến tính:} việc lọc theo bán kính hiện duyệt tuyến tính trên tập hồ sơ thay vì dùng chỉ mục địa lý, sẽ là điểm nghẽn khi số người dùng lớn.
  \item \textbf{Bộ nhớ đệm cục bộ theo tiến trình:} cache khoảng cách nằm trong bộ nhớ một tiến trình, chưa chia sẻ được khi chạy nhiều bản sao (scale-out).
\end{itemize}

\subsection{Bài học kinh nghiệm}
\label{subsection:6.1.4}

Quá trình thực hiện đồ án để lại nhiều bài học giá trị:
\begin{itemize}
  \item \textbf{Thiết kế cho điều kiện không lý tưởng:} người dùng thao tác đồng thời, dịch vụ ngoài có thể lỗi, quy trình có thể bị bỏ ngang --- buộc phải nghĩ về tương tranh, suy giảm an toàn và đối soát theo thời hạn ngay từ đầu thay vì xử lý chắp vá về sau.
  \item \textbf{Đơn giản nhưng đủ:} nhiều bài toán khó được giải bằng giải pháp nhẹ và đúng đắn (cập nhật nguyên tử thay cho khoá phân tán, bộ đếm phiên bản thay cho kho phiên) --- ``đủ dùng'' phù hợp quy mô đồ án hơn là công nghệ cồng kềnh.
  \item \textbf{Tính trung thực kỹ thuật:} phân biệt rõ phần đã hiện thực và phần còn là định hướng (ví dụ OTP/SMS) giúp báo cáo phản ánh đúng hiện trạng.
  \item \textbf{Năng lực full-stack:} phối hợp nhiều mảng (di động, máy chủ, cơ sở dữ liệu, thời gian thực, bản đồ, thông báo đẩy, đa ngôn ngữ) đòi hỏi tư duy hệ thống và quản lý độ phức tạp.
\end{itemize}

% ===========================================================================
\section{Hướng phát triển}
\label{section:6.2}

\subsection{Hoàn thiện các chức năng hiện có}
\label{subsection:6.2.1}

Trước mắt, các công việc giúp sản phẩm sẵn sàng cho vận hành thực tế gồm:
\begin{itemize}
  \item \textbf{Tích hợp cổng SMS} (ví dụ Twilio hoặc nhà cung cấp trong nước) để gửi OTP thật, hoàn tất luồng xác thực số điện thoại.
  \item \textbf{Phát hành bản iOS} và đưa ứng dụng lên cửa hàng (Google Play, App Store).
  \item \textbf{Chuyển lưu trữ tệp sang object storage/CDN} (ví dụ S3, Cloudinary) để tăng độ bền và tốc độ tải ảnh.
  \item \textbf{Bổ sung kiểm thử tự động và CI/CD} (unit, integration, kiểm thử đầu cuối) nhằm bảo đảm chất lượng khi mở rộng tính năng.
  \item \textbf{Tối ưu truy vấn vị trí} bằng chỉ mục địa lý \texttt{2dsphere} và truy vấn \texttt{\$nearSphere} để DB lọc trực tiếp thay vì quét tuyến tính.
  \item \textbf{Chuyển cache sang Redis} dùng chung và bổ sung Socket.IO adapter để có thể chạy nhiều bản sao máy chủ (scale-out) trong khi vẫn giữ thời gian thực.
\end{itemize}

\subsection{Các hướng nâng cấp và mở rộng}
\label{subsection:6.2.2}

Về lâu dài, sản phẩm có thể được nâng cấp theo các hướng:
\begin{itemize}
  \item \textbf{Theo dõi vị trí tình nguyện viên thời gian thực} trên bản đồ cho người nhận trong lúc giao hàng.
  \item \textbf{Ghép cặp \& định tuyến thông minh:} gợi ý đơn phù hợp, tối ưu lộ trình khi một tình nguyện viên nhận nhiều đơn.
  \item \textbf{Hệ thống uy tín \& an toàn thực phẩm:} huy hiệu/xếp hạng dựa trên đánh giá, nhắc hạn sử dụng, phân loại và cảnh báo dị ứng.
  \item \textbf{Bảng đo lường tác động xã hội:} thống kê số bữa ăn được chia sẻ, lượng thực phẩm cứu khỏi bị bỏ đi và lượng phát thải tránh được --- hiện thực hoá mục tiêu nêu ở Chương~1.
  \item \textbf{Tổng quát hoá nền tảng:} mở rộng sang quyên góp quần áo, sách vở, thuốc còn hạn\ldots dựa trên cùng mô hình ba bên đã chứng minh ở miền thực phẩm.
  \item \textbf{Hợp tác với tổ chức thiện nguyện/NGO} để mở rộng nguồn cung -- cầu và tăng độ tin cậy của nền tảng.
\end{itemize}

Tóm lại, đồ án không chỉ tạo ra một sản phẩm hoàn chỉnh giải quyết bài toán chia sẻ thực phẩm cho bối cảnh Việt Nam, mà còn đặt nền móng kiến trúc và các giải pháp kỹ thuật có thể tái sử dụng cho nhiều bài toán điều phối nguồn lực xã hội khác.

\ifSubfilesClassLoaded{\printbibliography{}}{}
\end{document}
